\section*{Hoofdstuk \ref{ch:g7:heart-and-circulation} --- Het hart en de bloedsomloop}

\begin{solution}{exo:g7:heart-and-circulation:1}
Rechterboezem: vangt het terugkerende \omterm{prop:g4:journey-of-food:absorb}{bloed} van het lichaam op.
Rechterkamer: pompt het naar de \omterm{def:g7:how-animals-breathe:four}{longen}. Linkerboezem: vangt het
zuurstofrijke \omterm{prop:g4:journey-of-food:absorb}{bloed} van de \omterm{def:g7:how-animals-breathe:four}{longen} op. Linkerkamer: pompt het het hele
lichaam rond --- de dikwandigste van de vier.
\end{solution}

\begin{solution}{exo:g7:heart-and-circulation:2}
Het zijn eenrichtingsdeuren die voorkomen dat het \omterm{prop:g4:journey-of-food:absorb}{bloed} terugstroomt. De
lub-dub zijn hun paarsgewijze sluitingen: de instroomkleppen die
dichtslaan als de \omterm{def:g7:heart-and-circulation:rooms}{kamers} drijven (lub), de uitstroomkleppen die dichtslaan
als ze ontspannen (dub).
\end{solution}

\begin{solution}{exo:g7:heart-and-circulation:3}
Longcirculatie: van de rechterkamer naar de \omterm{def:g7:heart-and-circulation:vessels}{haarvaten} van de \omterm{def:g7:how-animals-breathe:four}{longen}
(\omterm{def:g7:what-is-respiration:respiration}{zuurstof} laden, \omterm{def:g7:what-is-respiration:respiration}{koolstofdioxide} lossen), terug naar de linkerboezem.
Lichaamscirculatie: van de linkerkamer naar de \omterm{def:g7:heart-and-circulation:vessels}{haarvaten} van het hele
lichaam (afleveren, ophalen), terug naar de rechterboezem.
\end{solution}

\begin{solution}{exo:g7:heart-and-circulation:4}
\omterm{def:g7:heart-and-circulation:vessels}{Slagaders}: dikke, elastische wanden --- ze voeren het \omterm{prop:g4:journey-of-food:absorb}{bloed} onder druk van
het \omterm{def:g4:heart-and-blood:heart}{hart} weg. \omterm{def:g7:heart-and-circulation:vessels}{Haarvaten}: wanden van één \omterm{def:g5:first-look-at-cells:cell}{cel} dik --- de plaats van elke
uitwisseling. \omterm{def:g7:heart-and-circulation:vessels}{Aders}: wijd, met dunnere wanden en eenrichtingskleppen ---
ze brengen het \omterm{prop:g4:journey-of-food:absorb}{bloed} onder lage druk naar het \omterm{def:g4:heart-and-blood:heart}{hart} terug.
\end{solution}

\begin{solution}{exo:g7:heart-and-circulation:5}
Zij drijft de lichaamscirculatie aan --- de veel grotere lus, naar tenen
en \omterm{def:g5:brain-in-command:system}{hersenen} en terug --- dus moet zij het hardst duwen, en haar wand is
daarop gebouwd.
\end{solution}

\begin{solution}{exo:g7:heart-and-circulation:6}
In de \omterm{def:g7:heart-and-circulation:vessels}{haarvaten}. De \omterm{def:g7:heart-and-circulation:vessels}{haarvaten} van de \omterm{def:g7:how-animals-breathe:four}{longen}: \omterm{def:g7:what-is-respiration:respiration}{zuurstof} erin,
\omterm{def:g7:what-is-respiration:respiration}{koolstofdioxide} eruit; die van de \omterm{prop:g7:digestion-and-nutrients:absorption}{darmvlokken}: \omterm{def:g7:digestion-and-nutrients:families}{voedingsstoffen} erin; die
van de \omterm{def:g3:bones-and-muscles:muscle}{spieren}: brandstof en \omterm{def:g7:what-is-respiration:respiration}{zuurstof} eraf, \omterm{def:g7:what-is-respiration:respiration}{koolstofdioxide} erop (die van
de \omterm{def:g1:the-five-senses:sense}{huid}: warmte eruit).
\end{solution}

\begin{solution}{exo:g7:heart-and-circulation:7}
Rechterkamer --- de \omterm{def:g7:heart-and-circulation:vessels}{haarvaten} van de \omterm{def:g7:how-animals-breathe:four}{long} (\omterm{def:g7:what-is-respiration:respiration}{zuurstof} aan boord) ---
linkerboezem --- linkerkamer --- de halsslagaders --- de \omterm{def:g7:heart-and-circulation:vessels}{haarvaten} van de
\omterm{def:g5:brain-in-command:system}{hersenen} --- de halsaders --- rechterboezem. De uitwisseling van de
\omterm{def:g5:brain-in-command:system}{hersenen} neemt een groot, gestaag deel van de \omterm{def:g7:what-is-respiration:respiration}{zuurstof} en de \omterm{prop:g7:organs-at-work:bill}{glucose} en
laat \omterm{def:g7:what-is-respiration:respiration}{koolstofdioxide} achter --- de vaste opdracht van de veeleisende
commandant.
\end{solution}

\begin{solution}{exo:g7:heart-and-circulation:8}
De \omterm{prop:g4:heart-and-blood:pulse}{polsslag} is het elastische oprekken van de slagaderwand bij elke stoot
van de \omterm{def:g7:heart-and-circulation:rooms}{kamer} --- hoge druk die een rekbare dikke wand ontmoet. Bij de
\omterm{def:g7:heart-and-circulation:vessels}{aders} is de stoot verbruikt bij het oversteken van de \omterm{def:g7:heart-and-circulation:vessels}{haarvaten}: een lage,
vloeiende druk in een zachte wijde buis --- er is niets meer om te
kloppen.
\end{solution}

\begin{solution}{exo:g7:heart-and-circulation:9}
De klim omhoog: vanuit de \omterm{def:g1:my-body:parts}{benen} moet het aderbloed onder lage druk tegen
de zwaartekracht in klimmen, en de \omterm{def:g7:heart-and-circulation:rooms}{kleppen} houden elke gewonnen hoogte
vast. Bij een soldaat die stokstijf staat knijpen de kuitspieren de \omterm{def:g7:heart-and-circulation:vessels}{aders}
nooit, valt de klim stil, hoopt het \omterm{prop:g4:journey-of-food:absorb}{bloed} zich in de \omterm{def:g1:my-body:parts}{benen} op en zakt de
aanvoer naar de \omterm{def:g5:brain-in-command:system}{hersenen} --- de flauwte. Marcherende voeten zijn
aderpompen.
\end{solution}

\begin{solution}{exo:g7:heart-and-circulation:10}
\omterm{def:g7:heart-and-circulation:vessels}{Slagader} erin vanaf de lichaamscirculatie; erbinnen nemen de \omterm{def:g7:heart-and-circulation:vessels}{haarvaten}
van de \omterm{prop:g7:digestion-and-nutrients:absorption}{darmvlokken} de \omterm{def:g7:digestion-and-nutrients:families}{voedingsstoffen} van de maaltijd op
(\cref{prop:g7:digestion-and-nutrients:absorption}); de bijzondere
aansluiting: haar \omterm{def:g7:heart-and-circulation:vessels}{ader} loopt eerst naar de \omterm{prop:g7:digestion-and-nutrients:absorption}{lever} om te sorteren en te
bankieren, en pas daarna voegt het \omterm{prop:g4:journey-of-food:absorb}{bloed} zich bij het algemene retour naar
de rechterboezem. Het schema vertelt het verteringshoofdstuk opnieuw.
\end{solution}

\begin{solution}{exo:g7:heart-and-circulation:11}
Het \omterm{prop:g4:journey-of-food:absorb}{bloed} in de ruimtes stroomt er alleen doorheen; de eigen dikke wanden
van het \omterm{def:g4:heart-and-blood:heart}{hart}, die zonder rust werken, hebben net als elke \omterm{def:g3:bones-and-muscles:muscle}{spier} hun eigen
haarvatlevering nodig --- dus vertakken kleine voedingsslagaders zich over
het \omterm{def:g4:heart-and-blood:heart}{hart} zelf. Dichtgeslibd smoren die de eigen brandstofleiding van de
pomp: het enige orgaan dat nooit kan pauzeren is precies het orgaan dat
hun verstopping uithongert.
\end{solution}

\begin{solution}{exo:g7:heart-and-circulation:12}
Elke uitwisseling waar het lichaam van leeft --- \omterm{def:g7:what-is-respiration:respiration}{zuurstof} bij de
\omterm{def:g7:lungs-and-gas-exchange:alveoli}{longblaasjes}, \omterm{def:g7:digestion-and-nutrients:families}{voedingsstoffen} bij de \omterm{prop:g7:digestion-and-nutrients:absorption}{darmvlokken}, brandstof en afval bij
elke \omterm{def:g3:bones-and-muscles:muscle}{spier} --- gebeurt alleen door haarvatwanden heen; \omterm{def:g7:heart-and-circulation:vessels}{slagaders} en \omterm{def:g7:heart-and-circulation:vessels}{aders}
brengen het \omterm{prop:g4:journey-of-food:absorb}{bloed} slechts naar die wanden toe en weer weg, en het \omterm{def:g4:heart-and-blood:heart}{hart}
houdt het slechts in beweging. Pomp en snelwegen bestaan omwille van de
haarfijne rijstroken: leg het haarvatverkeer ergens stil en dat orgaan
gaat dood, hoe gaaf het leidingwerk ook is.
\end{solution}

\begin{solution}{pb:g7:heart-and-circulation:1}
\textbf{1.} De \omterm{prop:g4:heart-and-blood:pulse}{polsslag}: het tellen van de stoten in de \omterm{def:g7:heart-and-circulation:vessels}{slagaders} --- het
tempo van de pomp, aan een \omterm{def:g1:my-body:joint}{pols} afgelezen. De bloeddruk: de slagaderdruk
tijdens de stoot en ertussen --- de belasting van de snelwegen. De lub-dub
van de stethoscoop: de \omterm{def:g7:heart-and-circulation:rooms}{kleppen} die je hoort sluiten --- de deugdelijkheid
van de deuren.

\textbf{2.} Een lage rustpolsslag bij een fitte patiënt betekent een
sterke \omterm{def:g7:heart-and-circulation:rooms}{kamer} die per slag meer \omterm{prop:g4:journey-of-food:absorb}{bloed} verplaatst: het getrainde \omterm{def:g4:heart-and-blood:heart}{hart}
draait rustig stationair, met reserve achter de hand.

\textbf{3.} De deuren van de \omterm{def:g7:heart-and-circulation:rooms}{kleppen}: de instroomkleppen bij het drijven
van de \omterm{def:g7:heart-and-circulation:rooms}{kamers} (lub), de uitstroomkleppen bij hun ontspanning (dub) ---
vier deuren die in twee paren netjes dichtvallen.

\textbf{4.} Rechterboezem, rechterkamer, de \omterm{def:g7:heart-and-circulation:vessels}{haarvaten} van de \omterm{def:g7:how-animals-breathe:four}{longen},
linkerboezem, linkerkamer, de \omterm{def:g7:heart-and-circulation:vessels}{slagaders} van het lichaam, de \omterm{def:g7:heart-and-circulation:vessels}{haarvaten} van
het lichaam (bijvoorbeeld van een \omterm{def:g3:bones-and-muscles:muscle}{spier}), de \omterm{def:g7:heart-and-circulation:vessels}{aders} --- en thuis bij de
rechterboezem.

\textbf{5.} Elke slag verliest een deel van zijn stoot achterwaarts door
het lek, dus krimpen de leveringen van de lichaamscirculatie per slag;
het \omterm{def:g4:heart-and-blood:heart}{hart} compenseert met meer slagen en zwaarder werk, en het
uithoudingsvermogen zakt --- de bijgeruis is levering die als
terugstroom verloren gaat.

\textbf{6.} De bloeddruk omhoog (stijve, vernauwde snelwegen bieden
weerstand aan de stoot) en het oprekken van de \omterm{def:g7:heart-and-circulation:vessels}{slagader} bij de \omterm{prop:g4:heart-and-blood:pulse}{polsslag}
veranderd; de linkerkamer, die elke ronde tegen het verstijfde circuit in
duwt, draait overuren en wordt dikker.

\textbf{7.} Dichtslibben in de eigen voedingsslagaders van het \omterm{def:g4:heart-and-blood:heart}{hart}
smoort de brandstofleiding van de \omterm{def:g3:bones-and-muscles:muscle}{spier} die nooit kan rusten --- de
klassieke weg naar een hartinfarct.

\textbf{8.} De eenrichtingskleppen van de \omterm{def:g7:heart-and-circulation:vessels}{aders} en de ontbrekende
spierpomp van beweging: stilstaand knijpen de kuiten nooit, valt de
winst van de \omterm{def:g7:heart-and-circulation:rooms}{kleppen} stil, en lopen de enkels vol. Wandelen knijpt de
\omterm{def:g7:heart-and-circulation:vessels}{aders} ritmisch --- de verlichting is mechanisch.

\textbf{9.} Omdat de rondrit door het lichaam de \omterm{def:g7:what-is-respiration:respiration}{zuurstof} van het \omterm{prop:g4:journey-of-food:absorb}{bloed}
uitgeeft: alleen de longcirculatie kan haar bijladen. De \omterm{def:g7:how-animals-breathe:four}{longen}
overslaan zou verbruikt \omterm{prop:g4:journey-of-food:absorb}{bloed} met een lege bestelwagen op pad sturen ---
de dubbele lus is bijtanken-en-dan-bezorgen, afgedwongen door de twee
helften van het \omterm{def:g4:heart-and-blood:heart}{hart}.

\textbf{10.} \omterm{def:g7:heart-and-circulation:vessels}{Slagader}: dik, elastisch, hoge druk --- de uitgaande
snelwegen. \omterm{def:g7:heart-and-circulation:vessels}{Haarvat}: één \omterm{def:g5:first-look-at-cells:cell}{cel} dik, dat alles doorweeft --- de
bezorgstroken waar alle uitwisselingen gebeuren. \omterm{def:g7:heart-and-circulation:vessels}{Ader}: wijd, met
\omterm{def:g7:heart-and-circulation:rooms}{kleppen}, lage druk --- het retour.

\textbf{11.} Bij elk consult deed het probleem er precies daar toe waar
het de haarvatdienst raakte: het lek hongerde de leveringen uit, het
dichtslibben smoorde de aanvoerlijnen, het opeenhopen legde het retour
stil. Aflezingen, \omterm{def:g7:heart-and-circulation:rooms}{kleppen} en snelwegen zijn middelen; het
haarvatverkeer --- de uitwisselingen --- is het doel.

\textbf{12.} Beweeg dagelijks --- de pomp wordt sterker en de snelwegen
blijven soepel; eet evenwichtig --- de wanden van de \omterm{def:g7:heart-and-circulation:vessels}{slagaders} blijven
vrij van aanslag; geen rook --- de vaten vernauwen noch verstijven, en de
stoelen in het \omterm{prop:g4:journey-of-food:absorb}{bloed} blijven vrij voor \omterm{def:g7:what-is-respiration:respiration}{zuurstof}.
\end{solution}
